\section{Limitations and Future Work}
\label{sec:limitations}

We acknowledge several boundaries of the current study, each of which opens a concrete direction for future work.

\begin{itemize}
    \item \textit{Observational nature.} The RQ2 coherence premium is associational as unobserved confounders such as financial literacy, advisor channel, account constraints, and prior portfolio composition are not ruled out by the volatility, segment, or year controls. Similarly, the elevated 2018 mean discordance is hypothesised as a MiFID II transition effect but not directly tested against a pre-MiFID counterfactual. A natural next step would be to apply causal inference techniques such as instrumental variables or difference-in-differences on datasets that span a regulatory boundary.
    \item \textit{Metric and band-assignment design choices.} The tolerance $d \leq 1$ is set rather than derived, binarising a 4-point ordinal distance, and the asset-band assignment in Section~\ref{sec:setup:bandassignment} is a hierarchical heuristic. Future work could explore data-driven tolerance calibration and validate the band assignment against external MiFID II classifications from fund providers.
    \item \textit{Intervention scope and tuning.} This study extends only LightGCN, leaving sequential, transformer, and content-based recommenders as candidates for the same stratification and coherence-margin treatment. Within LightGCN, $\lambda = 1$ is one point on a Pareto frontier between PC and ranking-quality metrics, and the intervention inherits the RQ3 best-trial backbone hyperparameters rather than being re-tuned end-to-end. The reported gains are therefore a lower bound on what a full frontier sweep or end-to-end re-tuning could achieve.
    \item \textit{Generalisation and measurement scope.} Results are demonstrated on FAR-Trans under MiFID II's 4-band scale. Extending the evaluation to other jurisdictions, longer time horizons, or alternative asset universes would strengthen the generalisability of PC@$k$ as a suitability metric. Additionally, RQ2 measures realised return on executed Buys while RQ4 measures forward return on top-$k$ recommended slates that may never be acted on. Bridging this gap through online evaluation or counterfactual off-policy estimation is a promising direction.
\end{itemize}
